\section{Generalisation to $\F_2$-Linear Cell Operators}
\label{sec:generalisation}

The proof of \Cref{lem:equivariance} uses only two facts about the
operator $\bitop$:
\begin{itemize}[leftmargin=2em]
  \item $\bitop$ is $\F_2$-linear in the cell coordinate;
  \item $\bitop$ acts trivially on the row coordinate (i.e., the same
    map at every row).
\end{itemize}
The coefficient-permutation structure of \Cref{def:bitop} was nowhere
used. Hence the construction extends, verbatim, to any operator in the
$\F_2$-linear envelope of bit-ops.

\subsection{The $\F_2$-linear envelope}

\begin{definition}[$\F_2$-linear cell operator]
\label{def:f2-linear-op}
An $\F_2$-linear cell operator is any element of $\End_{\F_2}(\cellring)$.
\end{definition}

The space $\End_{\F_2}(\cellring) \cong \cellring \otimes_{\F_2}
\cellring^{\vee}$ has $\F_2$-dimension $32 \times 32 = 1024$, and is
spanned by the elementary maps $E_{j,k} : X^k \mapsto X^j$. In
particular it contains $\F_2$-linear combinations of bit-ops:

\begin{example}[SHA-256 $\sigma_0$]
\label{ex:sigma-0}
The SHA-256 small-sigma operator
\[
  \sigma_0(x) \;=\; \Rot_7(x) \,\oplus\, \Rot_{18}(x) \,\oplus\, \ShR_3(x)
\]
is the $\F_2$-linear combination $\Rot_7 + \Rot_{18} + \ShR_3$
in $\End_{\F_2}(\cellring)$. By linearity it lifts to a virtual column
\[
  \sigma_{0\,*} v \;=\; \Rot_{7\,*} v \,+\, \Rot_{18\,*} v \,+\, \ShR_{3\,*} v
\]
on cell-wise $\F_2$-arithmetic, and
\Cref{lem:equivariance} applies.
\end{example}

\subsection{Linearity of the construction}

\begin{lemma}[Linearity]
\label{lem:linearity}
Let $\bitop_1,\bitop_2 \in \End_{\F_2}(\cellring)$ and let $a_1, a_2
\in \F_2$. Define $\bitop := a_1 \bitop_1 + a_2 \bitop_2$. Then
\[
  \mle{\bitop_* v}(r) \;=\; a_1 \mle{\bitop_{1\,*} v}(r) \,+\, a_2 \mle{\bitop_{2\,*} v}(r),
\]
and consequently
\[
  \psia\!\big(\bitop(\lmle{v}(r))\big) \;=\; a_1\, \psia\!\big(\bitop_1(\lmle{v}(r))\big) \,+\, a_2\, \psia\!\big(\bitop_2(\lmle{v}(r))\big).
\]
\end{lemma}

\begin{proof}
Both equalities are immediate from $\F_2$-linearity of the MLE
evaluation map and of $\psia$.
\end{proof}

\Cref{lem:linearity} has a practical consequence: a UAIR that needs
the virtual column $\sigma_{0\,*} v$ does \emph{not} need to declare
three separate bit-op specs and combine them downstream. It can
instead declare a single \emph{compound} spec
$(j, \sigma_0)$ once $\End_{\F_2}(\cellring)$ is admitted as the
domain of \Cref{def:bitop}, with the verifier-side reconstruction
\[
  e \;=\; \psia\!\big(\sigma_0(\lmle{v_j}(r_0))\big)
        \;=\; \psia\!\big((\Rot_7 + \Rot_{18} + \ShR_3)(\lmle{v_j}(r_0))\big),
\]
still costing $O(32)$ $\F_q$-operations.

\subsection{Beyond $\F_2$-linearity}

The lemma fails for cell operators that are not $\F_2$-linear in the
cell coordinate. A bit-mask $x \mapsto x \wedge \chi$ for a fixed mask
$\chi \in \cellring$ is $\F_2$-linear and therefore covered by
\Cref{lem:linearity}. By contrast, a multiplicative operator such as
$x \mapsto x \cdot y$ in $\cellring$ for a witness-dependent
$y \in \cellring$ is bilinear, not $\F_2$-linear in $x$ alone, and
does not fit the construction directly.

The construction in this paper covers the $\F_2$-linear case, which
includes \emph{all} bit-permuting and bit-masking operators of
SHA-256 and SHA-3.

\section*{Conclusion}

We have described a virtual-column construction that exploits the
commutation between coefficient-position $\F_2$-linear endomorphisms
of $\cellring = \F_2[X]/(X^{32})$ and the multilinear-extension
operator over the row coordinate. The construction adds
\emph{(i)}~ one MLE source to the constraint-aggregation sumcheck and
to mp\_eval; \emph{(ii)}~ one field element of proof per spec;
\emph{(iii)}~ a closed-form verifier-side reconstruction at the final
evaluation point. It introduces no fresh commitment, no fresh
challenge, no extra protocol round, and no new soundness loss beyond
the standard per-source mp\_eval term. By
\Cref{lem:linearity}, it covers arbitrary $\F_2$-linear functions of
the $32$ coefficient bits of a cell, which encompasses all
bit-rotations, bit-shifts, bit-masks, and their XOR-combinations as
they appear in SHA-256 and related primitives.
